HealthGraphBench est un benchmark exploratoire consacré à trois tâches temporelles de santé réglementaire : nouveaux liens produit--problème dans MAUDE, déficiences G--L lors d'inspections de maisons de retraite suivies par CMS et nouveaux liens prescripteur--médicament dans Medicare Part D. Nous distinguons les résultats historiques du benchmark des nouvelles comparaisons réglées et d'analyses diagnostiques antérieures sur MAUDE.

Dans les résultats historiques MAUDE, la fréquence des voisins atteint un rappel micro@10 de 0,192045, contre 0,172840 pour GraphSAGE à trois époques. L'analyse historique de durée C sélectionne trente époques sur validation et obtient 0,210642 sur 2024--2025, contre 0,192045 pour les voisins ; l'écart de $+0,018597$ a un IC apparié conditionnel à 95\,\% de $[+0,011975;+0,025570]$. Cet IC ne se transfère pas aux nouveaux ajustements Q2. Dans cette comparaison, réglée sur validation, le BPR retenu atteint un rappel micro@10 moyen de 0,230551 en 2024 et de 0,239437 en 2025, devant les références fixes. Le GraphSAGE à agrégation moyenne atteint respectivement 0,204065 et 0,211825 ; le contrôle sans propagation atteint 0,193403 et 0,165388. Les valeurs Q2 des modèles stochastiques sont moyennées sur les trois graines, sans sélection de graine. Le diagnostic historique D1 conserve une perte presque plate ; les résultats Q2 n'établissent pas que cette quasi-platitude résultait d'un bug d'optimisation historique.

Sur CMS, dans l'analyse historique, l'ajout des agrégats de propriétaires augmentait l'AUC regroupée de 0,001318 (IC historique incluant zéro) ; les comparaisons Q2 montrent des écarts faibles et non uniformément positifs selon l'algorithme et l'année. Sur Part D, le BPR Q2 atteint 0,236182 au test 2024, contre 0,217860 pour la popularité par spécialité ; historiquement, cette heuristique dépassait le BPR (0,217860 contre 0,180789). Les familles réglables ont été sélectionnées sur validation ; les résultats restent exploratoires, car ces périodes avaient déjà été consultées. Aucun nouvel intervalle n'est calculé pour les ajustements Q2. Une prévision prospective distincte porte sur la publication future de relations Part D, non sur une issue clinique ; l'évaluation indépendante n'a pas encore été réalisée. Ces résultats ne démontrent ni supériorité générale des graphes ni utilité clinique.
